% =====================================================================
% TERRA -- product features by tier
%
% Build from this directory:
%   latexmk -pdf product_features_by_tier.tex
% latexmk lives in /Library/TeX/texbin, outside the default PATH.
%
% The same items as TODO.md at the repository root, ordered by tier. The
% complexity column is reported for each item and is NOT what orders the
% tables: section 2 states the criterion that does.
% =====================================================================

\documentclass[11pt,a4paper]{article}

\usepackage[utf8]{inputenc}
\usepackage[T1]{fontenc}
\usepackage[english]{babel}
\usepackage[margin=2.0cm]{geometry}
\usepackage{lmodern}
\usepackage{microtype}
\usepackage{array}
\usepackage{booktabs}
\usepackage{tabularx}
\usepackage{xltabular}
\usepackage{enumitem}
\usepackage{xcolor}
\usepackage[round]{natbib}
\usepackage[hidelinks]{hyperref}

\linespread{1.08}
\setlist{itemsep=2pt, topsep=4pt}

% RAGGED-RIGHT TABLE COLUMNS. The columns are narrow, and a justified column
% of that width opens gaps between words. \arraybackslash restores the row
% separator that \raggedright redefines.
\newcolumntype{L}[1]{>{\raggedright\arraybackslash}p{#1}}
\newcolumntype{Y}{>{\raggedright\arraybackslash}X}

% One tier table: rank, product, what it delivers, what it builds on, what it
% waits for, and its complexity. The complexity is the last column on purpose:
% it is information about each row, not the key the rows are sorted by.
%
% xltabular, not tabularx: a tabularx is one unbreakable box, so a table that
% did not fit under its section's paragraph moved whole to the next page and
% left the rest of the page empty. This one breaks between rows and repeats
% its header on the next page.
\newcommand{\tierhead}{%
  \toprule
  \textbf{\#} & \textbf{Product} & \textbf{Delivers} & \textbf{Builds on} &
  \textbf{Depends on} & \textbf{Complexity} \\
  \midrule}
\newenvironment{tiertable}{%
  \footnotesize
  \setlength{\LTpre}{8pt}\setlength{\LTpost}{8pt}
  \xltabular{\linewidth}{@{}L{0.4cm}L{2.4cm}YYL{2.3cm}L{1.9cm}@{}}
  \tierhead
  \endfirsthead
  \tierhead
  \endhead
  \bottomrule
  \endlastfoot
}{%
  \endxltabular
}

\title{TERRA Product Features by Tier}
\author{Jo\~ao Leonardi\textsuperscript{a}}
\date{\today}

\makeatletter
\newcommand{\affilnum}[1]{\textsuperscript{#1}}
\newcommand{\address}[1]{\gdef\@address{#1}}
\address{\affilnum{a}Computer Vision Research Center, iCEV Institute of Higher Education, Teresina, Brazil.}
\renewcommand{\maketitle}{%
  \begin{center}
    {\LARGE\bfseries\@title\par}
    \vspace{0.8em}
    {\large\@author\par}
    \vspace{0.3em}
    {\small\itshape\@address\par}
    \vspace{0.3em}
    {\small\@date\par}
  \end{center}
  \vspace{0.8em}}
\makeatother

\begin{document}

\maketitle

\begin{abstract}
\noindent
This document lists eleven candidate products for TERRA in three tiers. Tier~1
holds the three items that complete the current workflow (area, fields, jobs,
results) or were requested directly. Tier~2 holds the five items that the
product comparison brought to this project lists and TERRA lacks. Tier~3 holds
three extensions proposed during that analysis. Each item states what it
delivers, the code it builds on, what it depends on, and a qualitative
complexity (low, medium or high). The complexity is reported for every item and
does not determine the order. Within a tier, items are ordered by dependency.
\end{abstract}

\section{Current state}

\paragraph{Implemented.} Classification with per-pixel confidence; comparison
with MapBiomas; NDVI series with phenology (start, peak and end of season,
computed over a run's area); surface water; mineral map; field delineation;
the job queue (\texttt{internal/jobs}); and the single compositor graph for
every field of an area (\texttt{frontend/src/lib/fieldSets.ts}).

\paragraph{Not implemented.} Sentinel-1, Landsat or HLS imagery; error-adjusted
area estimation; training on user samples; management zones;
socio-environmental checks; vegetation anomaly. Surface water is the only
product that computes an anomaly.

\section{Tier criterion}

Tiers are assigned by where a request comes from and what it unblocks, not by
effort:

\begin{description}[leftmargin=1.6cm, style=nextline]
  \item[Tier 1] Completes the workflow already built (area $\rightarrow$
    fields $\rightarrow$ jobs $\rightarrow$ results), or was requested
    directly. Items in later tiers report through these.
  \item[Tier 2] Listed in the product comparison brought to this project and
    absent from TERRA.
  \item[Tier 3] Proposed during the analysis and listed in neither source.
\end{description}

\noindent Within a tier, an item precedes the items that depend on it. Where
no dependency separates two items, the order follows the reach of the result:
products usable across all fields before products for one use. The complexity
column is a judgement from reading the code; it states scope and gives no
schedule.

\section{Tier 1: completing the current workflow}

The per-field table comes first because every later per-field product needs a
place to report one row per field. The season dates are already computed in
every field job and need only that table. Vegetation health was requested
directly and is a new job kind that the queue and the single graph accept.

\begin{tiertable}
1 & Per-field table &
One row per field (area by class, season dates, mean confidence),
exportable to CSV. &
Per-field runs; the single compositor graph (\texttt{fieldSets.ts}). &
None. & Low \\
\addlinespace
2 & Season dates per field &
Emergence, peak and senescence for each field, with an uncertainty window. &
\texttt{phenology.py}: Savitzky--Golay smoothing and a 50\,\% amplitude
threshold, the rule TIMESAT applies \citep{jonsson2004}. It already runs in
every field job. &
Item 1 (display across fields). & Low to medium \\
\addlinespace
3 & Vegetation health &
NDVI/NDRE anomaly of each field against its own history and against
neighbouring fields of the same crop, as a map and a score per field. &
Sentinel-2 series and red-edge bands B05--B07; a new job kind
(\texttt{jobWork}, \texttt{JOB\_KINDS}). &
Item 1 (score per field); decision D1. & Medium \\
\end{tiertable}

\section{Tier 2: gaps against the product comparison}

Socio-environmental checks lead the tier because they depend on no model and
apply to every field. The error-adjusted area follows because its labelling
screen is the one the training on user samples reuses. Management zones use
only the field polygons and the series. Sentinel-1 closes the tier: it feeds
the season dates and vegetation health of Tier~1 under cloud and is not a
prerequisite for any other item in this tier.

\begin{tiertable}
1 & Socio-envi\-ron\-mental checks &
Overlap of each field, in hectares and with dates, with PRODES/DETER
(deforestation after 2008, Brazilian Forest Code; after 31~Dec~2020, EUDR),
CAR/SICAR (APP and legal reserve), IBAMA embargoes, indigenous lands and
conservation units. &
Field polygons; area geometry. Missing: per-area download of the public
layers (TerraBrasilis WFS) and the vector intersections. &
None. & Medium \\
\addlinespace
2 & Error-adjusted area &
Area of each class with a 95\,\% confidence interval, and user's and
producer's accuracy \citep{olofsson2014}. &
Classification runs. Missing: stratified random sample, a labelling screen
over the imagery, the estimator. The method is cited in
\texttt{mapbiomas.py} and not computed. &
None. & Medium to high \\
\addlinespace
3 & Management zones &
Three to five zones per field from several seasons of NDVI/EVI (fuzzy
k-means; \citealp{fridgen2004}), exportable to GeoJSON or shapefile for
variable-rate application. &
Field polygons; Sentinel-2 series. &
Multi-season series per field. & Medium \\
\addlinespace
4 & Training on user samples &
A classifier fitted to the region, with labelled fields (``soybean'',
``maize'') as training samples. &
Field polygons; the Random Forest path. Prithvi is used frozen. Missing:
labelling screen, retraining. &
Item 2 (labelling screen). & Medium (Random Forest) to high (network) \\
\addlinespace
5 & Sentinel-1 (radar) &
A VV/VH series independent of cloud: season cycle in the summer crop,
harvest detection from the drop in VH, water mapping under cloud. &
Planetary Computer RTC collection, served through the same STAC as
Sentinel-2. Missing: a sidecar slice. &
None. & High \\
\end{tiertable}

\section{Tier 3: extensions proposed in the analysis}

Event alerts read the series that vegetation health builds, so they follow it.
HLS densifies the series that the dates and the health products read. Yield
closes the list because it depends on data the user supplies.

\begin{tiertable}
1 & Field event alerts &
Abrupt drops in the series (hail, frost, lodging, early harvest) found by
break detection (BFAST; \citealp{verbesselt2010}). &
Per-field series; can be part of vegetation health. &
Tier~1 item 3. & Medium \\
\addlinespace
2 & HLS (Landsat + Sentinel-2) &
A revisit of two to three days, which narrows the uncertainty of season
dates and health in cloudy regions. &
NASA LP DAAC; the Earthdata token already used by the mineral map. &
Tier~1 items 2 and 3 (consumers). & Medium \\
\addlinespace
3 & Yield &
Calibrated only with the user's harvest data (yield monitor); without it, a
relative potential map. &
Per-field series and season dates. &
User harvest data. & High \\
\end{tiertable}

\section{Open decisions}

\begin{description}[leftmargin=1cm, style=nextline]
  \item[D1] Reference for the vegetation anomaly: the field's own history,
    neighbouring fields of the same crop, or both. The choice sets the data
    each health job reads.
\end{description}

\section{Limitations}

\begin{itemize}
  \item \textbf{Criterion.} Tiers reflect the state of the code, the direct
    requests and the product comparison available when this list was
    written. A new request or a change in scope reorders them.
  \item \textbf{Complexity.} Qualitative, from reading the code; no effort was
    measured, and no item was prototyped.
  \item \textbf{Field accuracy.} Every per-field product inherits the field
    delineation, whose accuracy over western Paran\'a is unmeasured: no
    hand-digitised reference fields exist for the region.
  \item \textbf{Season dates.} NDVI marks emergence and senescence, not sowing
    and harvest. Sowing precedes emergence by several days and harvest
    follows senescence; with a five-day revisit and cloud, the uncertainty is
    days to weeks.
  \item \textbf{Public layers.} PRODES, DETER, CAR and the embargo lists change
    on their own schedules. A check is valid for the date its layers were
    read, and has to state that date.
  \item \textbf{Yield.} Without the user's harvest data, no absolute yield
    figure is supported.
\end{itemize}

\begin{thebibliography}{9}
\providecommand{\natexlab}[1]{#1}

\bibitem[Fridgen et~al.(2004)]{fridgen2004}
Fridgen, J.~J., Kitchen, N.~R., Sudduth, K.~A., Drummond, S.~T., Wiebold,
W.~J. and Fraisse, C.~W. (2004).
Management Zone Analyst (MZA): software for subfield management zone
delineation.
\emph{Agronomy Journal}, 96(1), 100--108.

\bibitem[J\"onsson and Eklundh(2004)]{jonsson2004}
J\"onsson, P. and Eklundh, L. (2004).
TIMESAT: a program for analyzing time-series of satellite sensor data.
\emph{Computers \& Geosciences}, 30(8), 833--845.

\bibitem[Olofsson et~al.(2014)]{olofsson2014}
Olofsson, P., Foody, G.~M., Herold, M., Stehman, S.~V., Woodcock, C.~E. and
Wulder, M.~A. (2014).
Good practices for estimating area and assessing accuracy of land change.
\emph{Remote Sensing of Environment}, 148, 42--57.

\bibitem[Verbesselt et~al.(2010)]{verbesselt2010}
Verbesselt, J., Hyndman, R., Newnham, G. and Culvenor, D. (2010).
Detecting trend and seasonal changes in satellite image time series.
\emph{Remote Sensing of Environment}, 114(1), 106--115.

\end{thebibliography}

\end{document}
